\cotefile{157-4}
\scanpage{164}
\mspage{\unsure{38|39}}

d'autre part l'hypothèse sur $W$ implique
que $W$ est stable par facteurs directs, \deleted{\unsure{donc}}
$C$ l'est aussi\,; $\unsure{A|R}$ est clos à gauche, donc $\inserted{TC{=}}C\cap W$
aussi, donc par prop. 4.1. cor. 2 on a
$TC=\widetilde{TC}$, $TC$ est clos à gauche. On gagne
donc, par la dernière assertion du cor. 3 prop. \unsure{4.3|5.3|3.3}. (p. 35).

En fait, indépendamment des conditions
particulières sur $\unsure{\mathcal M|M},W$, on vient de prouver
ceci :

\uline{Lemme} 4.4.1 (Cor. de prop. 4.3) Soit $\mathcal M$ une
catégorie, $W$ un localiseur dans $\mathcal M$ \unsure{multiplicativement|modérément}
saturé et stable par facteurs directs. Alors
les structures de Quillen disons $(C,F)$
associées à $W$ s'obtiennent ainsi : on
prend une partie \deleted{\unsure{du} $\unsure{C|G}$ \unsure{g}} $C_0$ de $\mathrm{Fl}\mathcal M$
telle que
\[
\deleted{\unsure{(4.4.2)}}\qquad TF\overset{\mathrm{déf}}{=}C_{0*}\subset W,
\]
on pose
\[
C=\widetilde C_0=TF^*,\quad TC=C\cap W,\quad F=(TC)_{\unsure{*}},
\]
et on \deleted{vérifie que} doit se débrouiller
pour que $(C,TF)$ et $(TC,F)$ satisfassent
à la condition de factorisation de
Quillen, i.e. soient des pré-couples de Quillen.

[En fait $C,TF\inserted{\text{ et }F}$ étant clos \deleted{\unsure{par}} \unsure{par} construction,
et $TC=C\cap W$ étant stable par facteurs directs
($C$ et $W$ l'étant), ce sont \unsure{$=$|là} des \uline{couples} de
Quillen.]

% archive sheet number: 163
